\documentclass[12pt, a4paper]{exam}
% Corrigé rédigé pour le site Ressources NSI et SNT (CC BY-SA 4.0).
% Compilation : pdflatex (deux passes).
\usepackage[T1]{fontenc}
\usepackage[utf8]{inputenc}
\usepackage{lmodern}
\usepackage[french]{babel}
\usepackage[a4paper,top=2.3cm,bottom=2.8cm,left=2cm,right=2cm,headsep=0.6cm,footskip=1.2cm]{geometry}
\usepackage{xcolor}
\usepackage{listings}
\usepackage{graphicx}
\usepackage{array}
\usepackage{amsmath}
\usepackage{amssymb}
\usepackage{enumitem}
\usepackage{needspace}
\usepackage{tikz}
\usetikzlibrary{positioning,arrows.meta,calc}
\usepackage[most]{tcolorbox}

% Coloration du code : même style que les sujets
\definecolor{codegreen}{rgb}{0,0.6,0}
\definecolor{codegray}{rgb}{0.5,0.5,0.5}
\definecolor{codepurple}{rgb}{0.58,0,0.82}
\definecolor{backcolour}{rgb}{0.95,0.95,0.92}
\lstdefinestyle{mystyle}{
    backgroundcolor=\color{backcolour},
    commentstyle=\color{codegreen}\itshape,
    keywordstyle=\color{magenta}\bfseries,
    numberstyle=\tiny\color{codegray},
    stringstyle=\color{codepurple},
    basicstyle=\footnotesize\ttfamily,
    breakatwhitespace=false,
    breaklines=true,
    postbreak=\mbox{\textcolor{codegray}{$\hookrightarrow$}\space},
    captionpos=b,
    keepspaces=true,
    columns=fullflexible,
    numbers=left,
    numbersep=5pt,
    xleftmargin=1.5em,
    showspaces=false,
    showstringspaces=false,
    showtabs=false,
    tabsize=4,
    upquote=true,
    extendedchars=true,
    literate={à}{{\`a}}1 {â}{{\^a}}1 {é}{{\'e}}1 {è}{{\`e}}1 {ê}{{\^e}}1 {ë}{{\"e}}1
             {î}{{\^i}}1 {ï}{{\"i}}1 {ô}{{\^o}}1 {ù}{{\`u}}1 {û}{{\^u}}1 {ç}{{\c c}}1
             {É}{{\'E}}1 {À}{{\`A}}1 {œ}{{\oe}}1 {…}{{\ldots}}1
}
\lstset{style=mystyle, language=Python}
\newcommand{\code}[1]{\texttt{#1}}

% Encadrés du corrigé (lisibles en noir et blanc)
\tcbset{corrbase/.style={enhanced, breakable, sharp corners, left=2mm, right=2mm, top=1mm, bottom=1mm,
  fonttitle=\bfseries\small, coltitle=black, before skip=2mm, after skip=3mm}}
\newtcolorbox{reponse}[1][Réponse]{corrbase, unbreakable, colback=gray!6, colframe=black, colbacktitle=gray!22, boxrule=0.8pt, title={#1}}
\newtcolorbox{erreur}[1][Erreur fréquente]{corrbase, colback=white, colframe=black, boxrule=0.4pt, leftrule=3pt, title={#1}, colbacktitle=white, attach title to upper={\ :\ }}
\newtcolorbox{retenir}[1][À retenir]{corrbase, colback=white, colframe=black, boxrule=0.4pt, borderline={0.4pt}{1.2pt}{black}, title={#1}, colbacktitle=white, attach title to upper={\ :\ }}

% En-tête d'une question : numéro, points, rappel de la question
\newcommand{\Q}[3]{\par\needspace{9\baselineskip}\bigskip\noindent{\large\bfseries Question #1}\ifx&#2&\else\ \textnormal{(#2)}\fi\par\nopagebreak\smallskip\noindent\textit{#3}\par\nopagebreak\smallskip}
\newcommand{\Qp}[2]{\par\medskip\noindent\textbf{#1}\ \textit{#2}\par\nopagebreak\smallskip}
\newcommand{\methode}{\par\noindent\textbf{Méthode.}\ }
\newcommand{\application}{\par\noindent\textbf{Application.}\ }

% Arbres
\tikzset{arbre/.style={every node/.style={draw,circle,minimum size=7mm,inner sep=1pt,font=\small},level distance=11mm, thick},
  marque/.style={fill=gray!35, line width=1.6pt},
  titre/.style={draw=none, rectangle, inner sep=2pt}}

\raggedbottom
\setlength{\parindent}{0pt}
\setlength{\parskip}{3pt}
\setlist{itemsep=1pt, topsep=2pt}

\title{Corrigé --- CC Term NSI - Arbre Binaire de Recherche}
\author{Babash - NSI Tani Malandi et Johanne Bakalara}
\date{Février 2025 \textperiodcentered{} Durée du contrôle : 30 minutes \textperiodcentered{} 20 points}

\begin{document}
\pagestyle{foot}
\footer{}{\footnotesize Babash - NSI Tani Malandi et Johanne Bakalara \textperiodcentered{} CC BY-SA 4.0 \textperiodcentered{} page \thepage}{}
\maketitle
\thispagestyle{foot}
\begin{center}\footnotesize\itshape Sujet d'après le sujet de NSI du baccalauréat 2023 (jour 2, Polynésie, code 23-NSIJ2PO1). Corrigé original.\end{center}

\begin{retenir}[Comment lire ce corrigé]
Pour chaque question : la question rappelée en italique, la réponse attendue dans un cadre gris, puis l'explication (la méthode en étapes, puis son application aux trois arbres). Les cadres « Erreur fréquente » montrent les pièges classiques. Essaie d'abord de refaire la question seul·e, puis compare.
\end{retenir}

\section*{Les trois arbres du sujet}

\begin{center}
\begin{tikzpicture}[arbre, level 1/.style={sibling distance=34mm}, level 2/.style={sibling distance=17mm}, level 3/.style={sibling distance=9mm}]
\node {20}
  child {node {7} child {node {4} child {node {2}} child[missing]} child {node {15} child {node {11}} child[missing]}}
  child {node {53} child {node {18}} child {node {80} child[missing] child {node {83}}}};
\node[titre, font=\bfseries] at (0,0.8) {Arbre 1};
\end{tikzpicture}
\hfill
\begin{tikzpicture}[arbre, level 1/.style={sibling distance=30mm}, level 2/.style={sibling distance=15mm}, level 3/.style={sibling distance=9mm}]
\node {25}
  child {node {10} child {node {5}} child {node {15} child {node {7}} child {node {18}}}}
  child {node {40} child {node {30}} child {node {50}}};
\node[titre, font=\bfseries] at (0,0.8) {Arbre 2};
\end{tikzpicture}
\end{center}
\begin{center}
\begin{tikzpicture}[arbre, level 1/.style={sibling distance=40mm}, level 2/.style={sibling distance=20mm}, level 3/.style={sibling distance=10mm}]
\node {7}
  child {node {3} child {node {1}} child {node {5}}}
  child {node {15} child {node {12} child {node {9} child {node {8}} child[missing]} child {node {13}}} child {node {20} child[missing] child {node {25}}}};
\node[titre, font=\bfseries] at (0,0.8) {Arbre 3};
\end{tikzpicture}
\end{center}

\begin{retenir}[Le vocabulaire des arbres]
\begin{itemize}
\item Un \textbf{nœud} est un rond de l'arbre. Il porte une valeur, appelée aussi \textbf{clé}.
\item La \textbf{racine} est le nœud tout en haut : c'est le seul nœud qui n'a pas de parent.
\item Une \textbf{feuille} est un nœud qui n'a aucun enfant (ni à gauche, ni à droite).
\item La \textbf{taille} d'un arbre est le nombre total de ses nœuds.
\item La \textbf{profondeur} d'un nœud est le nombre d'arêtes (de traits) entre la racine et ce nœud. La racine a la profondeur 0.
\item La \textbf{hauteur} d'un arbre est la plus grande profondeur de ses nœuds : c'est le nombre d'arêtes du plus long chemin entre la racine et une feuille.
\end{itemize}
\textbf{Attention aux conventions.} Certains cours et certains sujets de bac comptent les nœuds au lieu des arêtes : la racine a alors la profondeur 1, et un arbre réduit à sa racine a la hauteur 1. Avec cette convention, on ajoute 1 à toutes les profondeurs et à toutes les hauteurs de ce corrigé. Au bac, la convention est toujours rappelée dans l'énoncé : lis-la avant de répondre.
\end{retenir}

\Q{1}{1 point}{Quelle est la racine des arbres 1, 2 et 3 ?}
\begin{reponse}
Arbre 1 : \textbf{20}. \qquad Arbre 2 : \textbf{25}. \qquad Arbre 3 : \textbf{7}.
\end{reponse}
La racine est le nœud placé tout en haut, celui dont partent toutes les branches. Elle n'a pas de parent : aucun trait ne monte au-dessus d'elle.

\Q{2}{1 point}{Quelles sont les feuilles des arbres 1, 2 et 3 ?}
\begin{reponse}
Arbre 1 : \textbf{2, 11, 18, 83}. \\
Arbre 2 : \textbf{5, 7, 18, 30, 50}. \\
Arbre 3 : \textbf{1, 5, 8, 13, 25}.
\end{reponse}
\methode
\begin{enumerate}
\item Parcours les nœuds un par un.
\item Pour chacun, regarde s'il part un trait vers le bas.
\item S'il n'en part aucun, c'est une feuille.
\end{enumerate}
\application Dans l'arbre 1, le nœud 4 n'est pas une feuille : il a un enfant, 2. De même, 15 a un enfant (11) et 80 a un enfant (83). Un nœud qui n'a qu'un seul enfant n'est donc pas une feuille.

\begin{erreur}
Prendre pour des feuilles tous les nœuds de la dernière ligne, ou seulement eux. Dans l'arbre 1, la feuille 18 est sur la troisième ligne, et 83 sur la quatrième : les feuilles ne sont pas toutes au même niveau.
\end{erreur}

\Q{3}{1 point}{Quelle est la taille des arbres 1, 2 et 3 ?}
\begin{reponse}
Arbre 1 : \textbf{10}. \qquad Arbre 2 : \textbf{9}. \qquad Arbre 3 : \textbf{11}.
\end{reponse}
La taille est le nombre de nœuds. Compte-les ligne par ligne pour ne pas en oublier :
\begin{itemize}
\item Arbre 1 : $1 + 2 + 4 + 3 = 10$ (20 ; 7, 53 ; 4, 15, 18, 80 ; 2, 11, 83).
\item Arbre 2 : $1 + 2 + 4 + 2 = 9$ (25 ; 10, 40 ; 5, 15, 30, 50 ; 7, 18).
\item Arbre 3 : $1 + 2 + 4 + 3 + 1 = 11$ (7 ; 3, 15 ; 1, 5, 12, 20 ; 9, 13, 25 ; 8).
\end{itemize}

\Q{4}{1 point}{Quelle est la profondeur du nœud 7 dans les arbres 1, 2 et 3 ?}
\begin{reponse}
Arbre 1 : \textbf{1}. \qquad Arbre 2 : \textbf{3}. \qquad Arbre 3 : \textbf{0}. \\
\small(Avec la convention « la racine a la profondeur 1 » : 2, 4 et 1.)
\end{reponse}
\methode
\begin{enumerate}
\item Pars de la racine et descends jusqu'au nœud 7.
\item Compte les traits (les arêtes) que tu franchis.
\end{enumerate}
\application
\begin{itemize}
\item Arbre 1 : 20 $\to$ 7. Un seul trait : profondeur 1.
\item Arbre 2 : 25 $\to$ 10 $\to$ 15 $\to$ 7. Trois traits : profondeur 3.
\item Arbre 3 : 7 est la racine. On ne descend pas : profondeur 0.
\end{itemize}

\Q{5}{1 point}{Quelle est la hauteur des arbres 1, 2 et 3 ?}
\begin{reponse}
Arbre 1 : \textbf{3}. \qquad Arbre 2 : \textbf{3}. \qquad Arbre 3 : \textbf{4}. \\
\small(Avec la convention « un arbre réduit à sa racine a la hauteur 1 » : 4, 4 et 5.)
\end{reponse}
\methode
\begin{enumerate}
\item Cherche la feuille la plus basse de l'arbre.
\item Compte les traits entre la racine et cette feuille : c'est la hauteur.
\end{enumerate}
\application
\begin{itemize}
\item Arbre 1 : le plus long chemin est 20 $\to$ 7 $\to$ 4 $\to$ 2 (ou 20 $\to$ 53 $\to$ 80 $\to$ 83) : 3 traits.
\item Arbre 2 : 25 $\to$ 10 $\to$ 15 $\to$ 7 : 3 traits.
\item Arbre 3 : 7 $\to$ 15 $\to$ 12 $\to$ 9 $\to$ 8 : 4 traits.
\end{itemize}

\Q{6}{3 points}{Dire pour chacun des arbres si c'est un arbre binaire de recherche. Justifier.}
\begin{reponse}
\begin{itemize}
\item Arbre 1 : \textbf{non}. Le nœud 18 est dans le sous-arbre droit de la racine 20, alors que $18 < 20$.
\item Arbre 2 : \textbf{non}. Le nœud 7 est dans le sous-arbre droit du nœud 10, alors que $7 < 10$.
\item Arbre 3 : \textbf{oui}. Pour chaque nœud, toutes les clés à sa gauche sont plus petites et toutes les clés à sa droite sont plus grandes.
\end{itemize}
\end{reponse}
\textbf{La définition.} Un arbre binaire de recherche (ABR) est un arbre binaire où, \textbf{pour chaque nœud} :
\begin{itemize}
\item toutes les clés de son sous-arbre gauche (pas seulement son enfant gauche) sont plus petites que sa clé ;
\item toutes les clés de son sous-arbre droit sont plus grandes que sa clé.
\end{itemize}
\methode
\begin{enumerate}
\item Prends chaque nœud, en commençant par la racine.
\item Regarde \textbf{tous} les nœuds situés sous lui à gauche : ils doivent être plus petits.
\item Regarde \textbf{tous} les nœuds situés sous lui à droite : ils doivent être plus grands.
\item Dès qu'un nœud ne respecte pas la règle, l'arbre n'est pas un ABR : cite ce nœud dans ta justification.
\end{enumerate}

\begin{center}
\begin{tikzpicture}[arbre, level 1/.style={sibling distance=30mm}, level 2/.style={sibling distance=15mm}, level 3/.style={sibling distance=8mm}]
\node[marque] {20}
  child {node {7} child {node {4} child {node {2}} child[missing]} child {node {15} child {node {11}} child[missing]}}
  child {node {53} child {node[marque] {18}} child {node {80} child[missing] child {node {83}}}};
\node[titre, font=\small, align=center] at (0,0.9) {Arbre 1 : 18 est à droite de 20};
\end{tikzpicture}
\hspace{1cm}
\begin{tikzpicture}[arbre, level 1/.style={sibling distance=28mm}, level 2/.style={sibling distance=14mm}, level 3/.style={sibling distance=8mm}]
\node {25}
  child {node[marque] {10} child {node {5}} child {node {15} child {node[marque] {7}} child {node {18}}}}
  child {node {40} child {node {30}} child {node {50}}};
\node[titre, font=\small, align=center] at (0,0.9) {Arbre 2 : 7 est à droite de 10};
\end{tikzpicture}
\end{center}

\application
\begin{itemize}
\item \textbf{Arbre 1.} Le nœud 53 a bien 18 à sa gauche ($18 < 53$) : localement, tout semble juste. Mais 18 est aussi dans le sous-arbre \emph{droit} de la racine 20, où toutes les clés doivent être plus grandes que 20. Or $18 < 20$ : la règle n'est pas respectée.
\item \textbf{Arbre 2.} Le nœud 15 a bien 7 à sa gauche ($7 < 15$). Mais 7 est aussi dans le sous-arbre droit de 10, où toutes les clés doivent être plus grandes que 10. Or $7 < 10$.
\item \textbf{Arbre 3.} On vérifie chaque nœud : à gauche de 7, il y a 1, 3, 5 (tous $< 7$) ; à droite, 8, 9, 12, 13, 15, 20, 25 (tous $> 7$). À gauche de 15 : 8, 9, 12, 13 ($< 15$) ; à droite : 20, 25 ($> 15$). À gauche de 12 : 8, 9 ; à droite : 13. Et ainsi de suite : tout est respecté.
\end{itemize}

\begin{retenir}[Astuce pour vérifier]
Un arbre binaire est un ABR \textbf{si et seulement si} son parcours infixe donne les clés dans l'ordre croissant. Tu peux t'en servir pour contrôler ta réponse avec la question 9 :
l'infixe de l'arbre 1 contient « 20, 18 » et celui de l'arbre 2 contient « 10, 7 » : ces deux suites ne sont pas croissantes. L'infixe de l'arbre 3 est croissant.
\end{retenir}

\begin{erreur}
Comparer chaque nœud seulement avec ses deux enfants. Dans les arbres 1 et 2, chaque parent est bien rangé par rapport à ses enfants directs, et pourtant ce ne sont pas des ABR : il faut comparer avec \textbf{tous} les nœuds du sous-arbre.
\end{erreur}

\Q{7}{3 points}{Donner pour chacun des arbres le parcours en largeur.}
\begin{reponse}
Arbre 1 : 20, 7, 53, 4, 15, 18, 80, 2, 11, 83 \\
Arbre 2 : 25, 10, 40, 5, 15, 30, 50, 7, 18 \\
Arbre 3 : 7, 3, 15, 1, 5, 12, 20, 9, 13, 25, 8
\end{reponse}
\methode Le parcours en largeur lit l'arbre \textbf{étage par étage}, comme un texte :
\begin{enumerate}
\item on écrit la racine (profondeur 0) ;
\item puis tous les nœuds de profondeur 1, de gauche à droite ;
\item puis tous ceux de profondeur 2, de gauche à droite, et ainsi de suite.
\end{enumerate}
\application Pour l'arbre 3 : étage 0 : 7 ; étage 1 : 3, 15 ; étage 2 : 1, 5, 12, 20 ; étage 3 : 9, 13, 25 ; étage 4 : 8.

\Q{8}{3 points}{Donner pour chacun des arbres le parcours en profondeur préfixe.}
\begin{reponse}
Arbre 1 : 20, 7, 4, 2, 15, 11, 53, 18, 80, 83 \\
Arbre 2 : 25, 10, 5, 15, 7, 18, 40, 30, 50 \\
Arbre 3 : 7, 3, 1, 5, 15, 12, 9, 8, 13, 20, 25
\end{reponse}
\textbf{Les trois parcours en profondeur.} Ils suivent tous la même règle, appliquée à chaque nœud ; seul change le moment où l'on écrit le nœud :
\begin{center}
\renewcommand{\arraystretch}{1.3}
\begin{tabular}{|l|l|}
\hline
\textbf{Préfixe} & nœud, puis sous-arbre gauche, puis sous-arbre droit \\ \hline
\textbf{Infixe} & sous-arbre gauche, puis nœud, puis sous-arbre droit \\ \hline
\textbf{Suffixe} (ou postfixe) & sous-arbre gauche, puis sous-arbre droit, puis nœud \\ \hline
\end{tabular}
\end{center}
Le préfixe (le mot « pré » veut dire « avant ») écrit le nœud \emph{avant} ses sous-arbres.

\begin{retenir}[Pour contrôler tes parcours]
\begin{itemize}
\item Chaque parcours contient \textbf{tous} les nœuds, une seule fois : compte-les, tu dois retrouver la taille (question 3).
\item Le préfixe \textbf{commence} toujours par la racine ; le suffixe \textbf{finit} toujours par la racine.
\end{itemize}
\end{retenir}

\textbf{Méthode du tour de l'arbre.} Dessine un trait qui fait le tour de l'arbre en partant à gauche de la racine et en longeant chaque nœud. En préfixe, tu écris un nœud quand tu passes à sa \textbf{gauche} ; en infixe, quand tu passes \textbf{dessous} ; en suffixe, quand tu passes à sa \textbf{droite}.

\application pas à pas pour l'arbre 2 :
\begin{enumerate}
\item On écrit la racine \textbf{25}, puis on traite tout le sous-arbre gauche (racine 10), puis tout le sous-arbre droit (racine 40).
\item Sous-arbre de 10 : on écrit \textbf{10}, puis son sous-arbre gauche (\textbf{5}), puis son sous-arbre droit (racine 15) : \textbf{15}, \textbf{7}, \textbf{18}.
\item Sous-arbre de 40 : \textbf{40}, \textbf{30}, \textbf{50}.
\end{enumerate}
On obtient : 25, 10, 5, 15, 7, 18, 40, 30, 50.

\Q{9}{3 points}{Donner pour chacun des arbres le parcours en profondeur infixe.}
\begin{reponse}
Arbre 1 : 2, 4, 7, 11, 15, 20, 18, 53, 80, 83 \\
Arbre 2 : 5, 10, 7, 15, 18, 25, 30, 40, 50 \\
Arbre 3 : 1, 3, 5, 7, 8, 9, 12, 13, 15, 20, 25
\end{reponse}
\methode Pour chaque nœud : d'abord tout son sous-arbre gauche, puis le nœud, puis tout son sous-arbre droit. On commence donc par descendre le plus à gauche possible.

\application pour l'arbre 1 :
\begin{enumerate}
\item On descend à gauche depuis 20 : 7, puis 4, puis 2. Le nœud 2 n'a pas d'enfant gauche : on écrit \textbf{2}.
\item On remonte : \textbf{4} (4 n'a pas d'enfant droit), puis \textbf{7}.
\item Sous-arbre droit de 7 (racine 15) : d'abord son gauche \textbf{11}, puis \textbf{15}.
\item Le sous-arbre gauche de 20 est fini : on écrit \textbf{20}.
\item Sous-arbre droit de 20 (racine 53) : \textbf{18}, \textbf{53}, puis le sous-arbre de 80 : \textbf{80} (pas d'enfant gauche), \textbf{83}.
\end{enumerate}

\begin{erreur}
Croire que le parcours infixe est « toujours dans l'ordre croissant ». C'est vrai seulement pour un ABR. Les arbres 1 et 2 n'en sont pas : leur infixe n'est pas croissant, et il ne faut surtout pas le « corriger » en le triant.
\end{erreur}

\Q{10}{3 points}{Donner pour chacun des arbres le parcours en profondeur suffixe (postfixe).}
\begin{reponse}
Arbre 1 : 2, 4, 11, 15, 7, 18, 83, 80, 53, 20 \\
Arbre 2 : 5, 7, 18, 15, 10, 30, 50, 40, 25 \\
Arbre 3 : 1, 5, 3, 8, 9, 13, 12, 25, 20, 15, 7
\end{reponse}
\methode Pour chaque nœud : d'abord tout son sous-arbre gauche, puis tout son sous-arbre droit, et le nœud en dernier. Un nœud n'est donc écrit que lorsque tous ses descendants ont été écrits.

\application pour l'arbre 3 :
\begin{enumerate}
\item Sous-arbre gauche de 7 (racine 3) : \textbf{1}, \textbf{5}, puis \textbf{3}.
\item Sous-arbre droit de 7 (racine 15). Son sous-arbre gauche (racine 12) : le sous-arbre de 9 donne \textbf{8}, \textbf{9} ; puis \textbf{13} ; puis \textbf{12}. Son sous-arbre droit (racine 20) : \textbf{25}, puis \textbf{20}. Puis \textbf{15}.
\item On termine par la racine : \textbf{7}.
\end{enumerate}


\end{document}
